{
Java is one of the most widely used and powerful programming languages in the field of computer science. Developed by Sun Microsystems and later acquired by Oracle Corporation, Java is a high-level, object-oriented programming language that follows the principle of “Write Once, Run Anywhere” (WORA). This feature is possible due to the Java Virtual Machine (JVM), which allows Java programs to run on any platform that supports the JVM, making it a platform-independent language. Over the years, Java has evolved into a robust tool for building a wide range of applications—from desktop software and web applications to mobile apps and enterprise-level systems.

The primary objective of this assignment is to explore and implement fundamental Java programming concepts such as classes and objects, inheritance, polymorphism, abstraction, exception handling, file operations, and data structures like arrays or collections. Through this assignment, students gain hands-on experience in writing clean, modular, and efficient code that adheres to object-oriented design principles. Java's extensive library support and built-in functionalities provide a solid foundation for learning best programming practices and improving problem-solving abilities.

In this assignment, a particular problem or real-life scenario has been modeled using appropriate classes and methods. Emphasis is placed on writing readable and maintainable code using comments, indentation, and proper naming conventions. The assignment also includes the algorithm or logic behind the implementation, input-output structure, and a discussion of how the program works. Additional features such as user interaction through console input, exception handling for robustness, and file input/output (if applicable) have been integrated to reflect practical software development techniques.

Overall, this Java assignment not only strengthens the theoretical understanding of programming concepts but also helps in developing critical thinking and debugging skills that are essential for any aspiring computer science professional. It lays the groundwork for advanced study in Java frameworks like Spring Boot, JavaFX, and Android development, thereby bridging the gap between academic learning and real-world software engineering.
}
